\section{Operational measurements and confidence ranking}
\subsection{Local quality and completion time}
The local performance experiment evaluates 11 strata, four checkpoints, three batch sizes and eight rounds on the same A100 80GB PCIe, yielding 1,056 cells and 98,304 measured decisions. Each stratum uses 64 fixed measured states and disjoint warm-up states. Model processes are fresh per block; model order is balanced, CPU affinity and thread counts are fixed, and an external observer records GPU, worker-core and SMT-sibling activity.

The timing boundary includes uncached serialization, tokenization, transfer, inference and native structured-answer construction, with GPU synchronization. It excludes model loading, disk output, validation, network transport and queueing. Batch-one p50/p95 characterize request completion; batch-eight and batch-32 work divided by prediction time characterize fixed-batch throughput. Round statistics are summarized by their median and an eight-round bootstrap interval. Repeated states increase timing replication, not the effective accuracy sample size.

Figure~\ref{fig:operational} compares quality and latency within six declared workloads. For BoolQ, median-of-round request p50 values are 10.30, 8.55, 36.86 and 33.51 milliseconds for Laya English, Laya Multilingual, Llama and Qwen respectively. The quality coordinate belongs to the same fixed workload. It is not appropriate to infer a cross-task winner from the horizontal position alone: systems can trade reference agreement for lower completion time, and that trade-off changes with the task.

The host is shared. No foreign GPU compute process or sustained CPU/SMT threshold violation was observed, but six isolated residual-CPU intervals exceeded 15\% without meeting the fixed two-interval rejection rule. These intervals remain included. All raw timing records, monitor samples and code/environment hashes are retained, and all 1,056 cells independently reconcile with the generated tables.

\subsection{Selective prediction within a source}
A software action can be deferred when the model is uncertain. Figure~\ref{fig:selective} shows empirical accepted-set error as coverage grows under maximum class probability. All tied probabilities enter together, avoiding artificial improvements from an arbitrary within-tie ordering. The curve is computed separately for each model and source; intervals at displayed landmarks condition on the selected threshold.

These curves distinguish two failure modes that accuracy alone conflates: frequent errors at all confidence levels and errors concentrated among uncertain cases. They support a workload-specific investigation of escalation policies. A threshold selected from these test curves would be a post-hoc decision rule, so a deployable policy requires a separate calibration/development split and a held-out check. Candidate-code probabilities are evaluated as observed ranking statistics, not assumed calibrated equivalents of native decision confidence.
